% ============ §11 DVRs and Dedekind rings (原书 pp.78-86) ============
\section{DVRs and Dedekind rings}\label{sec:11}\index{discrete valuation rings, section 11}\index{Dedekind rings, section 11}

A valuation ring whose value group is isomorphic to $\Z$ is called
a \emph{discrete valuation ring} (DVR).\index{discrete valuation ring} Discrete refers to the fact that the value
group is a discrete subgroup of $\R$, and has nothing to do with the
$m$-adic topology of the local ring being discrete.

\begin{theorem}\label{thm:11.1}
Let $R$ be a valuation ring. Then the following conditions are
equivalent.
\begin{enumerate}[(1)]
  \item $R$ is a DVR;
  \item $R$ is a principal ideal domain;
  \item $R$ is Noetherian.
\end{enumerate}
\end{theorem}

\begin{proof}
Let $K$ be the field of fractions of $R$ and $m$ its maximal ideal.

$(1) \Rightarrow (2)$ Let $v_R$ the additive valuation of $R$ having value group $\Z$; this is
called the \emph{normalised additive valuation} corresponding to $R$. There exists
$t \in R$ such that $v_R(t) = 1$. For $0 \neq x \in m$, the valuation $v_R(x)$ is a positive
integer, say $v_R(x) = n$; then $v_R(x/t^n) = 0$, so that we can write $x = t^nu$ with $u$ a
unit of $R$. In particular $m = tR$. Let $I \neq (0)$ be any ideal of $R$; then
$\{v_R(a) \mid 0 \neq a \in I\}$ is a set of non-negative integers, and so has a smallest
element, say $n$. If $n = 0$ then $I$ contains a unit of $R$, so that $I = R$. If $n > 0$ then
there exists an $x \in I$ such that $v_R(x) = n$; then $I = xR = t^nR$. Therefore $R$ is a
principal ideal domain, and moreover every non-zero ideal of $R$ is a power
of $m = tR$.

$(2) \Rightarrow (3)$ is obvious.

$(3) \Rightarrow (2)$ In general, given any two ideals of a valuation ring, one contains
the other, so that any finitely generated ideal $a_1R + \cdots + a_nR$ is equal to one
of the $a_iR$, and therefore principal. Hence, if $R$ is Noetherian every ideal of $R$
is principal.

$(2) \Rightarrow (1)$ We can write $m = xR$ for some $x$. Now if we set $I = \bigcap_{\nu \ge 1} x^\nu R$
then this is also a principal ideal, so that we can write $I = yR$. If we set
$y = xz$, then from $y \in x^\nu R$ we get $z \in x^{\nu-1}R$, and since this holds for every
$\nu$, we have $z \in I$, hence we can write $z = yu$. Since $y = xz = xyu$, we have $y(1
- xu) = 0$, but then since $x \in m$ we must have $y = 0$, and therefore $I = (0)$.
Because of this, for every non-zero element $a \in R$, there is a well-defined
integer $\nu \ge 0$ such that $a \in x^\nu R$ but $a \notin x^{\nu+1}R$; we then set $v(a) = \nu$. It is not
difficult to see that if $a, b, c, d \in R - \{0\}$ satisfy $a/b = c/d$ then
\[
  v(a) - v(b) = v(c) - v(d);
\]
therefore setting $v(\xi) = v(a) - v(b)$ for $\xi = a/b \in K^*$ gives a map $v\colon K^* \to \Z$
which can easily be seen to be an additive valuation of $K$ whose valuation
ring is $R$. The value group of $v$ is clearly $\Z$, so that $R$ is a DVR.
\end{proof}

If $R$ is a DVR with maximal ideal $m$ then an element $t \in R$ such that
$m = tR$ is called a \emph{uniformising element} of $R$.\index{uniformising element}

\begin{remark}
A valuation ring $S$ whose maximal ideal $m_S$ is principal does not
have to be a DVR. To obtain a counter-example, let $K$ be a field, and $R$
a DVR of $K$; set $k = R/m_R$, and suppose that $\bar{R}$ is a DVR of $k$. Now let
$S$ be the composite of $R$ and $\bar{R}$. Let $f$ be a uniformising element of $R$, and
$g \in S$ be any element mapping to a uniformising element $\bar{g}$ of $\bar{R}$. Then
$\eta = fR \subset m_R \subset S \subset R$, and $m_S/m_R = \bar{g}\bar{R} = \sqrt{\bar{g}S/\bar{m}_R}$, and so
\[
  m_S = m_R + gS.
\]
On the other hand $g^{-1} \in R$, so that for any $h \in m_R$ we have $h/g \in m_R \subset S$, and
hence $m_R \subset gS$, so that
\[
  m_S = gS.
\]
However, $m_R = fR$ is not finitely generated as an ideal of $S$, being generated
by $f, fg^{-1}, fg^{-2}, \ldots$. The value group of $S$ is $\Z^2$, with the valuation
$v\colon K^* \to \Z^2$ given by
\[
  v(x) = (n, m), \quad \text{where } n = v_R(x) \text{ and } m = v_{\bar{R}}(\bar{\varphi}(xf^{-n})),
\]
where $\varphi\colon R \to R/m_R = k$ is the natural homomorphism.
\end{remark}

The previous theorem gives a characterisation of DVRs among valuation
rings; now we consider characterisations among all rings.

\begin{theorem}\label{thm:11.2}
Let $R$ be a ring; then the following conditions are
equivalent:
\begin{enumerate}[(1)]
  \item $R$ is a DVR;
  \item $R$ is a local principal ideal domain, and not a field;
  \item $R$ is a Noetherian local ring, $\dim R > 0$ and the maximal ideal $m_R$ is
        principal;
  \item $R$ is a one-dimensional normal Noetherian local ring.
\end{enumerate}
\end{theorem}

\begin{proof}
We saw $(1) \Rightarrow (2)$ in the previous theorem; $(2) \Rightarrow (3)$ is obvious.

$(3) \Rightarrow (1)$ Let $xR$ be the maximal ideal of $R$. If $x$ were nilpotent then we
would have $\dim R = 0$, and hence $x^\nu \neq 0$ for all $\nu$. By the Krull intersection
theorem (Theorem~\ref{thm:8.10}, (i)) we have $\bigcap_{\nu \ge 1} x^\nu R = (0)$, so that for $0 \neq y \in R$
there is a well-determined $\nu$ such that $y \in x^\nu R$ and $y \notin x^{\nu+1}R$. If $y = x^\nu u$, then
since $u \notin xR$ it must be a unit. Similarly, for $0 \neq z \in R$ we have $z = x^\mu v$, with $v$ a
unit. Therefore $yz = x^{\mu+\nu}uv \neq 0$, and so $R$ is an integral domain. Finally,
any element $\xi$ of the fraction field of $R$ can be written $\xi = x^\nu u$, with $u$ a unit of
$R$ and $\nu \in \Z$, and it is easy to see that setting $v(\xi) = \nu$ defines an additive
valuation of the field of fractions of $R$ whose valuation ring is $R$.

$(1) \Rightarrow (4)$ In a DVR the only ideals are $(0)$ and the powers of the maximal
ideal, so that the only prime ideals of $R$ are $(0)$ and $m_R$, and hence
$\dim R = 1$. By the previous theorem $R$ is Noetherian, and it is normal
because it is a valuation ring.

$(4) \Rightarrow (3)$ By assumption $R$ is an integral domain. Write $K$ for the field of
fractions and $m$ for the maximal ideal of $R$. Then $m \neq 0$, so that by Theorem~\ref{thm:8.10}, $m \neq m^2$; choose some $x \in m - m^2$. Since $\dim R = 1$ the only prime
ideals of $R$ are $(0)$ and $m$, so that $m$ must be a prime divisor of $xR$, and there
exists $y \in R$ such that $xR:y = m$. Set $a = yx^{-1}$; then $a \notin R$, but $am \subset R$. Now
we set $m^{-1} = \{\xi \in K \mid \xi m \subset R\}$, so that $R \subset m^{-1}$, and $R \neq m^{-1}$ since $a \in m^{-1}$.
Consider the ideal $m^{-1}m$ of $R$; since $R \subset m^{-1}$ we have $m \subset m^{-1}m$. If we
had $m = m^{-1}m$ then we would get $am \subset m$, and then $a$ would be integral
over $R$ by Theorem~\ref{thm:2.1}, so that $a \in R$, which is a contradiction. Hence we
must have $m^{-1}m = R$. Moreover, $xm^{-1} \subset R$ is an ideal, and if $xm^{-1} \subset m$
then we would have $xR = xm^{-1}m \subset m^2$, contradicting $x \notin m^2$. Therefore
$xm^{-1} = R$, and hence $xR = xm^{-1}m = m$, so that $m$ is principal.
\end{proof}

Quite generally, if $R$ is an integral domain and $K$ its fields of fractions, we
say that an $R$-submodule $I$ of $K$ is a \emph{fractional ideal}\index{fractional ideal} of $R$ if $I \neq 0$, and
there exists a non-zero element $\alpha \in R$ such that $\alpha I \subset R$ (see Ex.~\ref{ex:3.4-ex}). As
an $R$-module we have $I \cong \alpha I$, so that if $R$ is a Noetherian integral domain
then any fractional ideal is finitely generated. For $I$ a fractional ideal we
set $I^{-1} = \{\alpha \in K \mid \alpha I \subset R\}$; we say that $I$ is \emph{invertible}\index{invertible ideal} if $I^{-1}I = R$.

\begin{theorem}\label{thm:11.3}
Let $R$ be an integral domain and $I$ a fractional ideal of $R$.
Then the following conditions are equivalent:
\begin{enumerate}[(1)]
  \item $I$ is invertible;
  \item $I$ is a projective $R$-module;
  \item $I$ is finitely generated, and for every maximal ideal $P$ of $R$, the
        fractional ideal $I_P = IR_P$ of $R_P$ is principal.
\end{enumerate}
\end{theorem}

\begin{proof}
$(1) \Rightarrow (2)$ If $I^{-1}I = R$ then there exist $a_i \in I$ and $b_i \in I^{-1}$ such that
$\sum_i a_ib_i = 1$. Then $a_1, \ldots, a_n$ generate $I$, since for any $x \in I$ we have $\sum (xb_i)a_i =
x$, and $xb_i \in R$. Let $F = Re_1 + \cdots + Re_n$ be the free $R$-module with basis
$e_1, \ldots, e_n$; we define the $R$-linear map $\varphi\colon F \to I$ by $\varphi(e_i) = a_i$, so that $\varphi$ is
surjective. Then we define $\psi\colon I \to F$ by writing $\psi_i\colon I \to R$ for the map
$\psi_i(x) = b_ix$, and setting $\psi(x) = \sum \psi_i(x)e_i$. We then have $\varphi\psi(x) = x$, so that $\varphi$
splits, and $I$ is isomorphic to a direct summand of the free module $F$, and
therefore projective.

$(2) \Rightarrow (1)$ Every $R$-linear map from $I$ to $R$ is given by multiplication by
some element of $K$ (prove this!). If we let $\varphi\colon F \to I$ be a surjective map from
a free module $F = \bigoplus Re_i$, by assumption there exists a splitting $\psi\colon I \to F$
such that $\varphi\psi = 1$. Write $\psi(x) = \sum \lambda_i(x)e_i$ for $x \in I$; then by what we have said,
each $\lambda_i$ determines a $b_i \in K$ such that $\lambda_i(x) = b_ix$, and since for each $x$ there
are only finitely many $i$ such that $\lambda_i(x) \neq 0$, we have $b_i = 0$ for all but finitely
many $i$. Letting $b_1, \ldots, b_n$ be the non-zero ones, we have $\sum a_ib_ix = x$ for all
$x \in I$, where $a_i = \varphi(e_i)$. Thus $\sum_i a_ib_i = 1$. Moreover, since $\lambda_iI = \lambda_i \subset R$ we
have $b_i \in I^{-1}$, and therefore $I^{-1}I = R$.

$(1) \Rightarrow (3)$ As we have already seen, $I$ is finitely generated. Now if $\sum a_ib_i = 1$
and $P$ is any prime ideal then at least one of $a_ib_i$ must be a unit of $R_P$,
and $I_P = a_iR_P$. Hence $I_P$ is a principal fractional ideal.

$(3) \Rightarrow (1)$ If $I$ is finitely generated then $(I^{-1})_P = (I_P)^{-1}$. Indeed, the
inclusion $\subset$ holds for any ideal; for $\supset$, if $I = a_1R + \cdots + a_nR$ and $x \in (I_P)^{-1}$
then $xa_i \in R_P$, so there exist $c_i \in R - P$ such that $xa_ic_i \in R$, so that setting
$c = c_1 \cdots c_n$ we have $(cx)a_i \in R$ for all $i$, which gives $cx \in I^{-1}$ and $x \in (I^{-1})_P$.
From the fact that $I_P$ is principal, we get $I_P \cdot (I_P)^{-1} = R_P$. Now if $II^{-1} \neq R$
then we can take a maximal ideal $P$ such that $II^{-1} \subset P$, and then
$I_P \cdot (I_P)^{-1} = I_P(I^{-1})_P \subset PR_P$, which is a contradiction. Thus we must have
$II^{-1} = R$.
\end{proof}

\begin{theorem}\label{thm:11.4}
Let $R$ be a Noetherian integral domain, and $P$ a non-zero
prime ideal of $R$. If $P$ is invertible then $\htop P = 1$ and $R_P$ is a DVR.
\end{theorem}

\begin{proof}
If $P$ is invertible the maximal ideal $PR_P$ of $R_P$ is principal, and
condition (3) of Theorem 2 is satisfied; thus $R_P$ is a DVR, and so $\dim R_P = 1$.
\end{proof}

\begin{theorem}\label{thm:11.5}
Let $R$ be a normal Noetherian domain. Then we have
\begin{enumerate}[(i)]
  \item all the prime divisors of a non-zero principal ideal have height $1$;
  \item $R = \bigcap_{\htop p = 1} R_p$.
\end{enumerate}
\end{theorem}

\begin{proof}
(i) Suppose $0 \neq a \in R$ and that $P$ is one of the prime divisors of $aR$;
then there exists an element $b \in R$ such that $aR:b = P$. We set $PR_P = m$, and
then $aR_P:b = m$, so that $ba^{-1} \notin R_P$.
If $ba^{-1}m \subset m$ then by the determinant trick $ba^{-1}$ is integral over $R_P$, which
contradicts the fact that $R_P$ is integrally closed. Thus $ba^{-1}m = R_P$, so that $m^{-1}m = R_P$, and
then by the previous theorem we get $\htop m = \htop P = 1$.

(ii) It is sufficient to prove that for $a, b \in R$ with $a \neq 0$, $b \in aR_P$ for every
height $1$ prime $P \in \Spec R$ implies $b \in aR$. Let $P_1, \ldots, P_r$ be the prime
divisors of $aR$, and let $aR = q_1 \cap \cdots \cap q_r$ be a primary decomposition of
$aR$, where $q_i$ is a $P_i$-primary ideal for each $i$. Then since $\htop P_i = 1$, we have
$b \in aR_{P_i} \cap R = q_i$ for $i = 1, \ldots, n$, and therefore $b \in \bigcap q_i = aR$.
\end{proof}

\begin{corollary*}
Let $R$ be a Noetherian domain. The following two conditions
are necessary and sufficient for $R$ to be normal:
\begin{enumerate}[(a)]
  \item for $P$ a height $1$ prime ideal, $R_P$ is a DVR;
  \item all the prime divisors of a non-zero principal ideal of $R$ have height $1$.
\end{enumerate}
\end{corollary*}

\begin{proof}
We have already seen necessity. For sufficiency, note that the proof
of (ii) above shows that (b) implies $R = \bigcap_{\htop p = 1} R_p$. Then by (a) each $R_p$
is normal, so that $R$ is normal.
\end{proof}

\begin{definition}
An integral domain for which every non-zero ideal is invertible
is called a \emph{Dedekind ring}\index{Dedekind ring} (sometimes \emph{Dedekind domain}).
\end{definition}

\begin{theorem}\label{thm:11.6}
For an integral domain $R$ the following conditions are
equivalent:
\begin{enumerate}[(1)]
  \item $R$ is a Dedekind ring;
  \item $R$ is either a field or a one-dimensional Noetherian normal domain;
  \item every non-zero ideal of $R$ can be written as a product of a finite
        number of prime ideals.
\end{enumerate}
Moreover, the factorisation into primes in (3) is unique.
\end{theorem}

\begin{proof}
$(1) \Rightarrow (2)$ Every non-zero ideal is invertible, and therefore finitely
generated, so that $R$ is Noetherian. Let $P$ be a non-zero prime ideal of $R$;
then by Theorem 4, the local ring $R_P$ is a DVR and $\htop P = 1$, and therefore
either $R$ is a field or $\dim R = 1$. Also, by Theorem~\ref{thm:4.7} we know that $R$ is the
intersection of the $R_P$ as $P$ runs through all the maximal ideals of $R$, but
since each $R_P$ is a DVR it follows that $R$ is normal.

$(2) \Rightarrow (1)$ If $R$ is a field there is no problem. If $R$ is not a field then for every
maximal ideal $P$ of $R$ the local ring $R_P$ is a one-dimensional Noetherian
local ring and is normal, so that by Theorem 2 it is a principal ideal ring.
Thus by Theorem 3, $R$ is a Dedekind ring.

$(1) \Rightarrow (3)$ Let $I$ be a non-zero ideal. If $I = R$ then we can view it as the
product of zero ideals; if $I$ is itself maximal then it is the product of just one
prime ideal. We have already seen that $R$ is Noetherian, so that we can use
descending induction on $l(I)$, that is assume that $I \neq R$ and that every ideal
strictly bigger than $I$ is a product of prime ideals. If $I \neq R$ then there is a
maximal ideal $P$ containing $I$, and $I \subset IP^{-1} \subset R$. If $IP^{-1} = I$ then using
$P^{-1}P = R$ we would have $I = IP$, and by NAK this would lead to a
contradiction. Hence $IP^{-1} \neq I$, so that by induction we can write $IP^{-1} =
Q_1 \cdots Q_s$, with $Q_i \in \Spec R$. Multiplying both sides by $P$ gives $I = Q_1 \cdots Q_sP$.

The proof of $(3) \Rightarrow (1)$ is a little harder, and we break it up into four steps.

\textbf{Step 1.} In general, any non-zero principal ideal $aR$ of an integral domain $R$
is obviously invertible. Moreover, suppose that $I$ and $J$ are non-zero
fractional ideals and $B = IJ$; then obviously $I$ and $J$ invertible implies $B$
invertible, but the converse also holds. To see this, from $I^{-1}J^{-1}B \subset R$ we
get $I^{-1}J^{-1} \subset B^{-1}$, and also from $B^{-1}IJ \subset R$ we get $B^{-1}I \subset J^{-1}$ and
$B^{-1}J \subset I^{-1}$; now if $B$ is invertible then multiplying the last two inclusions
together we get $B^{-1} = B^{-1}B^{-1}IJ \subset I^{-1}J^{-1}$, and hence $B^{-1} = I^{-1}J^{-1}$.
Therefore
\[
  R = BB^{-1} = IJI^{-1}J^{-1} = (II^{-1})(JJ^{-1}),
\]
and we must have $II^{-1} = JJ^{-1} = R$.

\textbf{Step 2.} Let $P$ be a non-zero prime ideal. Let us prove that if $I$ is an ideal
strictly bigger than $P$ then $IP = P$. For this it is sufficient to show that if
$I = P + aR$ with $a \notin P$ then $P \subset IP$. Consider expressions of $P^2$ and $a^2R + P$
as product of prime ideals, $P^2 = P_1 \cdots P_r$ and $a^2R + P = Q_1 \cdots Q_s$. Then $P_i$
and $Q_j$ are prime ideals containing $P^2$, and so are prime ideals strictly bigger than $P$,
We now set $\bar{R} = R/P$, and write $-$ to denote the image in $\bar{R}$ of
elements or ideals of $R$. Then we have
\[
  (*) \quad \bar{P}_1 \cdots \bar{P}_r = \bar{a}^2\bar{R} = \bar{Q}_1 \cdots \bar{Q}_s,
\]
and applying Step 1 to the domain $\bar{R}$ we find that $\bar{P}_i$ and $\bar{Q}_j$ are all invertible,
and are prime ideals of $\bar{R}$. We can suppose that $\bar{P}_1$ is a minimal element of
the set $\{\bar{P}_1, \ldots, \bar{P}_r\}$. Moreover, at least one of $\bar{Q}_1, \ldots, \bar{Q}_s$ is contained
in $\bar{P}_1$, so that we can assume that $\bar{Q}_1 \subset \bar{P}_1$, and, on the other hand, since
$\bar{Q}_1$ is also a prime ideal and $\bar{P}_1 \cdots \bar{P}_r \subset \bar{Q}_1$ we must have $\bar{Q}_1 \supset \bar{P}_i$
for some $i$. Then $\bar{P}_i \subset \bar{Q}_1 \subset \bar{P}_1$, and by the minimality of $\bar{P}_1$ we have $\bar{P}_i =
\bar{P}_1 = \bar{Q}_1$. Multiplying through both sides of $(*)$ by $\bar{P}_1^{-1}$ gives
\[
  \bar{P}_2 \cdots \bar{P}_r = \bar{Q}_2 \cdots \bar{Q}_s;
\]
proceeding in the same way, we see that $r = s$, and that after reordering the
$\bar{Q}_j$ we can assume that $\bar{P}_i = \bar{Q}_i$ for $i = 1, \ldots, r$. From this we get $P_i = Q_i$, and
$a^2R + P = (P + aR)^2 = P^2 + aP + a^2R$. Thus any element $x \in P$ can be
written
\[
  x = y + az + a^2t \quad \text{with } y \in P^2, z \in P \text{ and } t \in R.
\]
Since $a \notin P$ we must have $t \in P$, and then as required we have
$P \subset P^2 + aP = (P + aR)P$.

\textbf{Step 3.} Let $b \in R$ be a non-zero element; then in the factorisation
$bR = P_1 \cdots P_r$, every $P_i$ is a maximal ideal of $R$. Indeed, if $I$ is any ideal
strictly greater than $P_i$ then $IP_i = P_i$, and by Step 1 $P_i$ is invertible, so that
$I = R$.

\textbf{Step 4.} Let $P$ be a prime ideal of $R$, and $0 \neq a \in P$. If $aR = P_1 \cdots P_r$ with
$P_i \in \Spec R$ then $P$ must contain one of the $P_i$, but from Step 3 we know that
$P_i$ is maximal, so that $P = P_i$. We deduce that $P$ is a maximal ideal and is
invertible. If every non-zero prime ideal is invertible then any non-zero ideal
of $R$ is invertible, since it can be written as a product of primes. This
completes the proof of $(3) \Rightarrow (1)$.

Finally, if (1), (2) and (3) hold, then as we have seen in Step 2 above, the
uniqueness of factorisation into primes is a consequence of the fact that
the prime ideals of $R$ are invertible.
\end{proof}

\begin{theorem}[the Krull--Akizuki theorem]\label{thm:11.7}
Let $A$ be a one-dimensional
Noetherian integral domain with field of fractions $K$, let $L$ be a finite
algebraic extension field of $K$, and $B$ a ring with $A \subset B \subset L$; then $B$ is a
Noetherian ring of dimension at most $1$, and if $J$ is a non-zero ideal of $B$ then
$B/J$ is an $A$-module of finite length.
\end{theorem}

\begin{proof}
We follow the method of proof of Akizuki [1] in the linear algebra
formulation of [BS]. First of all we prove the following lemma.

\begin{lemma}\label{lem:11.1}
Let $A$ and $K$ be as in the theorem, and let $M$ be a torsion-free
$A$-module (see Ex.~\ref{ex:10.2-ex}) of rank $r < \infty$. Then for $0 \neq a \in A$ we have
\[
  l(M/aM) \le r \cdot l(A/aA).
\]
\end{lemma}

\begin{remark}
The \emph{rank} of a module $M$ over an integral domain $A$ is the maximal
number of elements of $M$ linearly independent over $A$; this is equal to the
dimension of the $K$-vector space $M \otimes_A K$.
\end{remark}

\begin{proof}[Proof of the lemma]
First we assume that $M$ is finitely generated. Choose
elements $\xi_1, \ldots, \xi_r \in M$ linearly independent over $A$ and set $E = \sum A\xi_i$; then
for any $\eta \in M$ there exists $t \in A$ with $t \neq 0$ such that $t\eta \in E$. If we set $C = M/E$
then from the assumption on $M$ we see that $C$ is also finitely generated, so
that $tC = 0$ for suitable $0 \neq t \in A$. Applying Theorem~\ref{thm:6.4} to $C$, we can find
$C = C_0 \supset C_1 \supset \cdots \supset C_m = 0$, such that $C_{i-1}/C_i \cong A/p_i$, with $p_i \in \Spec A$.
Now $t \in p_i$, and since $A$ is one-dimensional each $p_i$ is maximal, so that
$l(C) = m < \infty$. If $0 \neq a \in A$ then the exact sequence
\[
  E/a^nE \to M/a^nM \to C/a^nC \to 0
\]
gives
\[
  (*) \quad l(M/a^nM) \le l(E/a^nE) + l(C) \quad \text{for all } n > 0.
\]
Now $E$ and $M$ are both torsion-free $A$-modules, and one sees easily that
$aM/a^{i+1}M \cong M/aM$, and similarly for $E$. Hence $(*)$ can be written
$n \cdot l(M/aM) \le n \cdot l(E/aE) + l(C)$ for all $n > 0$, which gives $l(M/aM) \le l(E/aE)$.
Since $E \cong A^r$ we have $l(E/aE) = r \cdot l(A/aA)$. This completes the proof in the
case that $M$ is finitely generated. If $M$ is not finitely generated, take a finitely
generated submodule $\bar{M}_0 = A\bar{\xi}_1 + \cdots + A\bar{\xi}_r$ of $\bar{M} = M/aM$. Then choos-
ing an inverse image $\omega_i$ in $M$ for each $\bar{\xi}_i$, and setting $M_0 = \sum A\omega_i$, we get
\[
  l(\bar{M}_0) = l(M_0/M_0 \cap aM) \le l(M_0/aM_0) \le r \cdot l(A/aA).
\]
The right-hand side is now independent of $\bar{M}_0$, so that $\bar{M}$ is in fact finitely
generated, and $l(\bar{M}) \le r \cdot l(A/aA)$.
\end{proof}

We return to the proof of the theorem. We can replace $L$ by the field of
fractions of $B$. Set $[L:K] = r$; then $B$ is a torsion-free $A$-module of rank $r$.
Hence by the lemma, for any $0 \neq a \in A$ we have $l_A(B/aB) < \infty$. Now if $J \neq 0$ is
an ideal of $B$ and $0 \neq b \in J$ then since $b$ is algebraic over $A$ it satisfies a
relation
\[
  a_0b^m + a_1b^{m-1} + \cdots + a_{m-1}b + a_m = 0 \quad \text{with } a_i \in A.
\]
$B$ is an integral domain, so that we can assume $a_0 \neq 0$. Then $0 \neq a_0 \in J \cap A$
and so
\[
  l_A(B/J) \le l_A(B/a_0B) < \infty.
\]
Moreover, one sees from $l_A(J/a_0B) \le l_A(B/a_0B) < \infty$ that $J/a_0B$
is a finite $B$-module; hence, $J$ itself is a finite $B$-module, and therefore $B$ is
Noetherian. If $P$ is a non-zero prime ideal of $B$ then $B/P$ is an Artinian
ring and an integral domain, and therefore a field. Thus $P$ is maximal and
$\dim B = 1$.
\end{proof}

\begin{corollary*}
Let $A$ be a one-dimensional Noetherian integral domain, $K$ its
field of fractions, and $L$ a finite algebraic extension field of $K$; write $B$ for the
integral closure of $A$ in $L$. Then $B$ is a Dedekind ring, and for any maximal
ideal $P$ of $A$ there are just a finite number of primes of $B$ lying over $P$.
\end{corollary*}

\begin{proof}
By the theorem $B$ is a one-dimensional Noetherian integral domain
and is normal by construction; hence it is a Dedekind ring. It is easy to see
that if we factorise $PB$ as a product $PB = Q_1^{e_1} \cdots Q_r^{e_r}$ of a finite number of
prime ideals, then $Q_1, \ldots, Q_r$ are all the prime ideals of $B$ lying over $P$.
\end{proof}

% ============ 习题 ============
\section*{Exercises to \S\ref{sec:11}}
\emph{Prove the following propositions.}

\begin{exercise}\label{ex:11.1-ex}
Let $A$ be a DVR, $K$ its field of fractions, and $\bar{K}$ an algebraic closure of $K$;
then any valuation ring of $\bar{K}$ dominating $A$ is a one-dimensional non-
discrete valuation ring.
\end{exercise}

\begin{exercise}\label{ex:11.2-ex}
Let $A$ be a DVR, $K$ its field of fractions, and $L$ a finite extension field of $K$;
then a valuation ring of $L$ dominating $A$ is a DVR.
\end{exercise}

\begin{exercise}\label{ex:11.3-ex}
Let $A$ be a DVR and $m$ its maximal ideal; then the $m$-adic completion $\hat{A}$ of
$A$ is again a DVR.
\end{exercise}

\begin{exercise}\label{ex:11.4-ex}
Let $v\colon K \to \R \cup \{\infty\}$ be an Archimedean additive valuation of a field $K$,
and let $c$ be a real number with $0 < c < 1$. For $\alpha, \beta \in K$, set $d(\alpha, \beta) = c^{v(\alpha-\beta)}$;
then $d$ satisfies the axioms for a metric on $K$ (that is $d(\alpha, \beta) \ge 0$, $d(\alpha, \beta)
= 0 \iff \alpha = \beta$, $d(\alpha, \beta) = d(\beta, \alpha)$ and $d(\alpha, \gamma) \le d(\alpha, \beta) + d(\beta, \gamma)$), and the to-
pology of $K$ defined by $d$ does not depend on the choice of $c$. Let $R$ be the
valuation ring of $v$ and $m$ its valuation ideal; if $R$ is a DVR then the
topology determined by $d$ restricts to the $m$-adic topology on $R$.
\end{exercise}

\begin{exercise}\label{ex:11.5-ex}
Any ideal in a Dedekind ring can be generated by at most two elements.
\end{exercise}

\begin{exercise}\label{ex:11.6-ex}
Let $A$ be the integral closure of $\Z$ in $\Q(\sqrt{10})$; then $A$ is a Dedekind ring but
not a principal ideal ring.
\end{exercise}

\begin{exercise}\label{ex:11.7-ex}
If a Dedekind ring $A$ is semilocal then it is a principal ideal ring.
\end{exercise}

\begin{exercise}\label{ex:11.8-ex}
A module over a Dedekind ring is flat if and only if it is torsion-free.
\end{exercise}

\begin{exercise}\label{ex:11.9-ex}
Let $A$ be an integral domain (not necessarily Noetherian). The following
two conditions are equivalent:
\begin{enumerate}[(1)]
  \item $A_P$ is a valuation ring for every maximal ideal $P$ of $A$;
  \item an $A$-module is flat if and only if it is torsion-free. (An integral
        domain satisfying these conditions is called a \emph{Pr\"ufer domain}.)
\end{enumerate}
\end{exercise}

\begin{exercise}\label{ex:11.10-ex}
A finite torsion-free module over a Dedekind ring is projective, and is
isomorphic to a direct sum of ideals.
\end{exercise}
